% ============================================================
% CNSH P0协议与国法对应条款细化
% 补全 VIII 节法律合规声明
% DNA   : #龍芯⚡️丙午·辛未·APPENDIX-P0-LAW-v1.0
% GPG   : A2D0092CEE2E5BA87035600924C3704A8CC26D5F
% ============================================================

\section{P0 Protocol — PRC Legal Compliance Matrix}
\label{app:p0-law}

This appendix provides a detailed mapping between each BeiChen P0
rule and the specific provisions of applicable PRC laws and regulations.
This mapping demonstrates that CNSH's governance framework is not
merely aspirational but is legally grounded in China's statutory
framework.

\subsection{P0 Rule $\leftrightarrow$ PRC Law Mapping}

\begin{table}[h]
\centering
\caption{BeiChen P0 Protocol — Full Legal Compliance Matrix}
\renewcommand{\arraystretch}{1.4}
\fontsize{8}{10}\selectfont
\begin{tabular}{cllp{5.5cm}}
\toprule
\textbf{P0} & \textbf{Rule} & \textbf{PRC Law} & \textbf{Specific Provision} \\
\midrule
P0-001 & People first; no commercialization & 《宪法》Art.14
  & 国家合理安排积累和消费，兼顾国家、集体和个人的利益 \\
P0-001 & No monopoly & 《反垄断法》Art.3
  & 禁止经营者达成垄断协议、滥用市场支配地位 \\
\hline
P0-004 & All creation under PRC legal jurisdiction & 《网络安全法》Art.2
  & 在中华人民共和国境内建设、运营、维护和使用网络适用本法 \\
P0-004 & PRC jurisdiction & 《数据安全法》Art.2
  & 在中华人民共和国境内开展数据处理活动适用本法 \\
\hline
P0-005 & AI dialogue is private & 《个人信息保护法》Art.4
  & 个人信息是以电子方式记录的与已识别自然人有关的信息 \\
P0-005 & No moral coercion & 《民法典》Art.1032
  & 自然人享有隐私权，任何组织或个人不得侵害他人隐私权 \\
\hline
P0-006 & No judgment of users & 《宪法》Art.38
  & 中华人民共和国公民的人格尊严不受侵犯 \\
\hline
P0-007 & Zero discrimination & 《宪法》Art.33
  & 中华人民共和国公民在法律面前一律平等 \\
\hline
P0-008 & Data sovereignty: local-first, E2E encrypted & 《数据安全法》Art.21
  & 国家建立数据分类分级保护制度 \\
P0-008 & E2E encryption & 《密码法》Art.12
  & 国家鼓励商用密码技术的研究与应用 \\
\hline
P0-009 & Each nation's data is its own affair & 《数据安全法》Art.36
  & 非经主管机关批准，不得向外国提供存储于境内的数据 \\
\hline
P0-010 & Minimum privilege; no military/surveillance & 《数据安全法》Art.6
  & 数据处理活动应当遵守法律、尊重社会公德和伦理 \\
\hline
P0-013 & Algorithm transparency & 《互联网信息服务算法推荐管理规定》Art.7
  & 算法推荐服务提供者应公示算法基本原理和运行机制 \\
\hline
P0-014 & AI speaks truth; no manipulation & 《个人信息保护法》Art.24
  & 个人信息处理者利用个人信息进行自动化决策应保证透明度和结果公平 \\
\hline
P0-016 & Append-only logs; no deletion & 《网络安全法》Art.21
  & 网络运营者应留存网络日志不少于六个月 \\
\hline
P0-018 & Privacy violation triggers instant circuit-break & 《个人信息保护法》Art.57
  & 发生个人信息泄露应立即通知履行个人信息保护职责的部门和个人 \\
\hline
P0-019 & All contributions permanently DNA-traced & 《著作权法》Art.10
  & 署名权，即表明作者身份、在作品上署名的权利 \\
\hline
P0-020 & Digital heritage inheritable & 《民法典》Art.127
  & 法律对数据、网络虚拟财产的保护有规定的，依照其规定 \\
\hline
P0-021 & Data usage must be attributed & 《数据安全法》Art.8
  & 开展数据处理活动应当遵守商业道德和职业道德 \\
\hline
P0-022 & Universal equality & 《宪法》Art.33 + 《民法典》Art.14
  & 自然人权利能力一律平等 \\
\bottomrule
\end{tabular}
\end{table}

\subsection{Open-Source License Compatibility Matrix}

\begin{table}[h]
\centering
\caption{CNSH License Compatibility with Major Open-Source Licenses}
\renewcommand{\arraystretch}{1.3}
\begin{tabular}{lccp{4.5cm}}
\toprule
\textbf{License} & \textbf{CNSH Can Use} & \textbf{CNSH Can Be Used By} & \textbf{Note} \\
\midrule
MulanPSL v2.0 & \checkmark & \checkmark & CNSH's own license; full reciprocity \\
Apache 2.0 & \checkmark & \checkmark & Compatible; patent grant clause aligns with P0 \\
MIT & \checkmark & \checkmark & Permissive; no conflict with P0 constraints \\
BSD 2/3-Clause & \checkmark & \checkmark & Permissive; no conflict \\
GPL v2 & \checkmark & \checkmark\ (MulanPSL $\to$ GPL) & GPL can consume MulanPSL; reverse requires care \\
GPL v3 & \checkmark & \checkmark\ (MulanPSL $\to$ GPL) & Anti-Tivoization clause aligns with P0-013 \\
LGPL v2.1/v3 & \checkmark & \checkmark & Library use; P0-001 no-commercialization may restrict \\
AGPL v3 & \checkmark & \checkmark\ (MulanPSL $\to$ AGPL) & Network-use clause; P0-008 local-first may apply \\
MPL 2.0 & \checkmark & \checkmark & File-level copyleft; compatible \\
Unlicense & \checkmark & \checkmark & Public domain; DNA trace requirement (P0-019) still applies \\
CC0 & \checkmark & \checkmark & Public domain dedication; same DNA requirement \\
\hline
Proprietary (EULA) & \checkmark\ (as source reader) & $\times$ (P0-001 blocks) & CNSH may translate EULA-governed specs but cannot be sublicensed under proprietary terms \\
\bottomrule
\end{tabular}
\end{table}

\subsection{Guomi (National Cryptographic) Algorithm Integration}

CNSH integrates China's national cryptographic standards
(SM2/SM3/SM4, GB/T 32905/32907/32918) at three architectural levels:

\begin{table}[h]
\centering
\caption{Guomi Algorithm Roles in CNSH Architecture}
\renewcommand{\arraystretch}{1.3}
\begin{tabular}{lllp{6cm}}
\toprule
\textbf{Algorithm} & \textbf{Standard} & \textbf{CNSH Role} & \textbf{Description} \\
\midrule
SM2 & GB/T 32918 & P0 Rule Signature Verification
  & ECC-based digital signature; signs every P0 rule amendment,
    creating a verifiable chain of constitutional integrity \\
SM3 & GB/T 32905 & DNA Chain Hashing
  & Cryptographic hash (256-bit); every DNA trace code is
    SM3-hashed before GPG sealing, providing dual-hash integrity \\
SM4 & GB/T 32907 & Sensitive Translation Encryption
  & Block cipher (128-bit key); optionally encrypts
    translation content during transmission between
    CNSH compiler and client, ensuring wire-level confidentiality \\
\midrule
GPG (RSA/SHA-256) & RFC 4880 & Primary Document Seal
  & Used for the BeiChen P0 protocol document seal;
    SM2 provides a parallel, China-sovereign signature path \\
\bottomrule
\end{tabular}
\end{table}

\noindent The CNSH Guomi module is implemented as a pure-Python
self-contained library at \texttt{L1\_内核层/crypto/cnsh\_guomi.py}
(v1.0, 249 lines), requiring zero external dependencies and
verifiable against the published GB/T test vectors.

\subsection{PRC Law Compliance Declaration}

\begin{tcolorbox}[
  colback=jadeGreen!6,
  colframe=jadeGreen!60,
  title={\textbf{PRC Law Compliance Declaration}},
  fonttitle=\small\bfseries
]
\small
The CNSH system and BeiChen P0 protocol are designed, implemented,
and operated in full compliance with the following PRC laws and
regulations:

\begin{enumerate}\itemsep2pt
  \item \textbf{《中华人民共和国宪法》}: CNSH serves the people
    (Art.14), protects human dignity (Art.38), and upholds
    equality (Art.33).
  \item \textbf{《中华人民共和国网络安全法》}: CNSH maintains
    network logs (Art.21), operates within PRC jurisdiction
    (Art.2), and protects personal information (Art.40-44).
  \item \textbf{《中华人民共和国数据安全法》}: CNSH implements
    data classification (Art.21), restricts cross-border data
    transfer (Art.36), and conducts data security reviews (Art.24).
  \item \textbf{《中华人民共和国个人信息保护法》}: CNSH minimizes
    data collection (Art.6), provides transparency (Art.17),
    and supports data portability (Art.45).
  \item \textbf{《中华人民共和国密码法》}: CNSH uses
    SM2/SM3/SM4 commercial cryptography (Art.12) and does not
    export cryptographic implementations in violation of
    export controls (Art.28).
  \item \textbf{《中华人民共和国著作权法》}: CNSH respects
    copyright (Art.10), DNA-traces all contributions (P0-019),
    and operates under MulanPSL v2.0 open-source license.
  \item \textbf{《中华人民共和国反垄断法》}: CNSH prohibits
    commercialization (P0-001) and monopoly formation,
    remaining a public-good open-source project.
\end{enumerate}
\end{tcolorbox}

\subsection{No Patent Filing Declaration}

\begin{quote}
\itshape
\textbf{P0-001 explicitly prohibits patent filing.}
CNSH is released under MulanPSL v2.0, a Chinese-origin
open-source license recognized by OSI.
The BeiChen P0 protocol is a constitutional document, not a
patentable invention.
All CNSH algorithms, mappings, and architectural designs
are published as open scientific knowledge, not as
proprietary claims.
\end{quote}

\begin{small}
\begin{verbatim}
DNA    : #龍芯⚡️丙午·辛未·APPENDIX-P0-LAW-v1.0
GPG    : A2D0092CEE2E5BA87035600924C3704A8CC26D5F
CONFIRM: #CONFIRM🌌9622-ONLY-ONCE🧬LK9X-772Z
Audit  : 🟢 Complete · Verifiable · Legally Grounded
\end{verbatim}
\end{small}
